% =============================================================
% preamble.tex — Преамбула монографии (СТУ 7.5–07–2021 СФУ)
% =============================================================

% --- Язык и шрифты ---
\usepackage[utf8]{inputenc}
\usepackage[T2A]{fontenc}
\usepackage[russian]{babel}
\usepackage{tempora}

% --- Геометрия страницы (СТУ: 30/10/20/20 мм) ---
\usepackage{geometry}
\geometry{
  left=30mm, right=10mm,
  top=20mm, bottom=20mm,
  nohead,
  footskip=10mm  % номер страницы — посередине нижнего поля
}

\raggedbottom

% topskip = зазор TeX'а между верхом текстового блока и baseline 1-й строки.
% По умолчанию для 14pt extbook = 14.5pt (≈5 мм), из-за чего заголовки глав
% (СОДЕРЖАНИЕ, ВВЕДЕНИЕ, 1 …) смещаются вниз от рамки.
% При малом topskip зазор исчезает; число строк на странице не меняется,
% т.к. для 14pt шрифта высота строки (~10pt) всегда > topskip.
\topskip=0.1pt

% --- Интервалы, абзац ---
\usepackage{setspace}
\singlespacing

\usepackage{indentfirst}
\setlength{\parindent}{1.25cm}
\setlength{\parskip}{0pt}

\tolerance=1000
\hbadness=1000
\emergencystretch=1.5em

% --- Математика ---
\usepackage[fleqn]{amsmath}
\usepackage{amssymb,amsthm}
\usepackage{mathtools}
\usepackage{upgreek}

% Прямые греческие буквы по умолчанию
\renewcommand{\alpha}{\upalpha}
\renewcommand{\beta}{\upbeta}
\renewcommand{\gamma}{\upgamma}
\renewcommand{\delta}{\updelta}
\renewcommand{\epsilon}{\upepsilon}
\renewcommand{\varepsilon}{\upvarepsilon}
\renewcommand{\zeta}{\upzeta}
\renewcommand{\eta}{\upeta}
\renewcommand{\theta}{\uptheta}
\renewcommand{\vartheta}{\upvartheta}
\renewcommand{\iota}{\upiota}
\renewcommand{\kappa}{\upkappa}
\renewcommand{\lambda}{\uplambda}
\renewcommand{\mu}{\upmu}
\renewcommand{\nu}{\upnu}
\renewcommand{\xi}{\upxi}
\renewcommand{\pi}{\uppi}
\renewcommand{\rho}{\uprho}
\renewcommand{\sigma}{\upsigma}
\renewcommand{\tau}{\uptau}
\renewcommand{\upsilon}{\upupsilon}
\renewcommand{\phi}{\upphi}
\renewcommand{\varphi}{\upvarphi}
\renewcommand{\chi}{\upchi}
\renewcommand{\psi}{\uppsi}
\renewcommand{\omega}{\upomega}
\renewcommand{\Gamma}{\Upgamma}
\renewcommand{\Delta}{\Updelta}
\renewcommand{\Theta}{\Uptheta}
\renewcommand{\Lambda}{\Uplambda}
\renewcommand{\Xi}{\Upxi}
\renewcommand{\Pi}{\Uppi}
\renewcommand{\Sigma}{\Upsigma}
\renewcommand{\Phi}{\Upphi}
\renewcommand{\Psi}{\Uppsi}
\renewcommand{\Omega}{\Upomega}

\setlength{\mathindent}{1.25cm}

% --- Формулы: автоматические отступы по СТУ ---
\makeatletter
\AtBeginDocument{%
  % Номинально 1 строка; minus-компонента позволяет TeX'у сжать
  % отступ до ~½ строки, чтобы формула не вылезала за нижнюю границу
  \abovedisplayskip=1\baselineskip plus 0pt minus 8pt\relax
  \abovedisplayshortskip=1\baselineskip plus 0pt minus 8pt\relax
  \belowdisplayskip=1\baselineskip plus 0pt minus 8pt\relax
  \belowdisplayshortskip=1\baselineskip plus 0pt minus 8pt\relax

  % Формула в начале страницы: БЕЗ отступа сверху
  \everydisplay\expandafter{%
    \the\everydisplay
    \ifdim\pagetotal<1pt\relax
      \abovedisplayskip=0pt\relax
      \abovedisplayshortskip=0pt\relax
    \fi
  }%
}
\makeatother

% Разрешить разрыв страницы перед формулой (по умолчанию 10000 = запрещено)
\predisplaypenalty=0

% Защита от висячих строк (вдов/сирот)
\widowpenalty=150
\clubpenalty=150

% --- Блок «где» после формул (пояснение переменных) ---
% Использование:
%   \begin{conditions}
%   $n$ -- порядковый номер элемента;\\
%   $r_n$ -- расстояние от центра, м.
%   \end{conditions}
\newlength{\wherewd}
\AtBeginDocument{\settowidth{\wherewd}{где\enspace}}

\newenvironment{conditions}{%
  \par\noindent где\enspace\ignorespaces
  \begin{minipage}[t]{\dimexpr\linewidth-\wherewd\relax}%
}{%
  \end{minipage}\par\noindent\ignorespaces
}

% --- Перечисления ---
\usepackage{enumitem}
\setlist[itemize]{
  label={-},
  leftmargin=0pt,
  itemindent=\dimexpr\parindent+\labelsep+0.5em\relax,
  labelsep=0.5em,
  listparindent=0pt,
  topsep=0pt,
  partopsep=0pt,
  parsep=0pt,
  itemsep=2pt,
}

\setlist[enumerate]{
  leftmargin=0pt,
  itemindent=\dimexpr\parindent+\labelsep+1.5em\relax,
  labelsep=0.5em,
  labelwidth=1.5em,
  listparindent=0pt,
  topsep=0pt,
  partopsep=0pt,
  parsep=0pt,
  itemsep=0pt,
}

% --- Контроль размещения плавающих объектов ---
\usepackage{placeins}          % \FloatBarrier

% --- Графика ---
\usepackage{graphicx}
\graphicspath{{figures/}}
\usepackage{tikz}
\usetikzlibrary{arrows.meta,calc,patterns}

% --- Таблицы ---
\usepackage{booktabs}
\usepackage{longtable}
\usepackage{array}
\usepackage{tabularx}

% --- Подписи ---
\usepackage{caption}
\captionsetup{labelsep=endash, font={normalsize}}
\captionsetup[figure]{name=Рисунок, position=bottom, justification=centering}
\captionsetup[table]{name=Таблица, position=top, justification=raggedright, singlelinecheck=false, margin=0pt}

% --- Нумерация по главам ---
\usepackage{chngcntr}
\counterwithin{equation}{chapter}
\counterwithin{figure}{chapter}
\counterwithin{table}{chapter}

% --- Содержание ---
\usepackage[titles]{tocloft}
\renewcommand{\contentsname}{СОДЕРЖАНИЕ}

\setlength{\cftbeforetoctitleskip}{0pt}
\setlength{\cftaftertoctitleskip}{\baselineskip}

\setlength{\cftbeforechapskip}{0pt}

\setlength{\cftchapindent}{0pt}
\setlength{\cftsecindent}{5mm}
\setlength{\cftsubsecindent}{10mm}

\renewcommand{\cftchapleader}{\cftdotfill{\cftdotsep}}
\renewcommand{\cftsecleader}{\cftdotfill{\cftdotsep}}
\renewcommand{\cftsubsecleader}{\cftdotfill{\cftdotsep}}

\renewcommand{\cftchapfont}{}
\renewcommand{\cftchappagefont}{}
\renewcommand{\cftsecfont}{}
\renewcommand{\cftsecpagefont}{}
\renewcommand{\cftsubsecfont}{}
\renewcommand{\cftsubsecpagefont}{}

% --- Заголовки ---
\usepackage{titlesec}

\renewcommand{\chaptername}{}
\renewcommand{\thechapter}{\arabic{chapter}}

\titleformat{\chapter}[block]
  {\normalfont\normalsize\bfseries}
  {\hspace*{1.25cm}\thechapter}{0.5em}{}
\titlespacing*{\chapter}{0pt}{0pt}{\baselineskip}

\titleformat{name=\chapter,numberless}[block]
  {\normalfont\normalsize\bfseries\centering}
  {}{0pt}{\MakeUppercase}
\titlespacing*{name=\chapter,numberless}{0pt}{0pt}{\baselineskip}

\titleformat{\section}[block]
  {\normalfont\normalsize\bfseries}
  {\hspace*{1.25cm}\thesection}{0.5em}{}
\titleformat{name=\section,numberless}[block]
  {\normalfont\normalsize\bfseries}
  {}{0pt}{\hspace*{1.25cm}}
\titlespacing*{\section}{0pt}{\baselineskip}{\baselineskip}

\titleformat{\subsection}[block]
  {\normalfont\normalsize\bfseries}
  {\hspace*{1.25cm}\thesubsection}{0.5em}{}
\titlespacing*{\subsection}{0pt}{\baselineskip}{\baselineskip}

% --- Библиография ---
\usepackage[
  backend=biber,
  style=gost-numeric,
  natbib=true,
  sorting=none
]{biblatex}
\addbibresource{bibliography.bib}

% Запятая вместо точки с запятой между номерами: [10, 11]
\DeclareDelimFormat{multicitedelim}{\addcomma\space}
% Убрать вертикальные отступы между записями
\setlength{\bibitemsep}{0pt}
\setlength{\bibparsep}{0pt}
% Убрать висячий отступ: все строки записи без доп. отступа
\defbibenvironment{bibliography}
  {\list{\printtext[labelnumberwidth]{\printfield{labelprefix}\printfield{labelnumber}.}}
     {\setlength{\labelwidth}{\labelnumberwidth}%
      \setlength{\leftmargin}{\labelwidth}%
      \setlength{\labelsep}{\biblabelsep}%
      \addtolength{\leftmargin}{\labelsep}%
      \setlength{\itemindent}{0pt}%
      \setlength{\itemsep}{\bibitemsep}%
      \setlength{\parsep}{\bibparsep}}%
    \renewcommand*{\makelabel}[1]{\hss##1}}
  {\endlist}
  {\item}
% Диапазоны страниц: дефис (134-144), разделители полей: короткое тире (--)
\AtBeginBibliography{%
  \renewcommand{\bibrangedash}{-}%
  \def\textemdash{\textendash}%
}

% --- Гиперссылки ---
\usepackage[hidelinks]{hyperref}

% --- Нумерация страниц ---
\usepackage{fancyhdr}
\pagestyle{fancy}
\fancyhf{}
\fancyfoot[C]{\thepage}
\renewcommand{\headrulewidth}{0pt}
\renewcommand{\footrulewidth}{0pt}

\fancypagestyle{plain}{%
  \fancyhf{}%
  \fancyfoot[C]{\thepage}%
  \renewcommand{\headrulewidth}{0pt}%
  \renewcommand{\footrulewidth}{0pt}%
}

% --- Float-страницы ---
\makeatletter
\setlength{\@fptop}{0pt}
\setlength{\@fpsep}{\baselineskip}
\setlength{\@fpbot}{0pt}
\makeatother

